\chapter{DAG --- Layer-Skipping Edges Silently Occluded By Intermediate Nodes}\label{record-289}

\noindent{\footnotesize\color{BookMuted}\path{docs/method/retro/DAG-layer-skipping-edges-occluded.md}\par}\smallskip

\begin{RecordMeta}
title: DAG — Layer-Skipping Edges Silently Occluded By Intermediate Nodes
lane: retro

\end{RecordMeta}

\section{Disposition}\label{record-289-disposition}

Fixed on \texttt{r\allowbreak{}e\allowbreak{}l\allowbreak{}e\allowbreak{}a\allowbreak{}s\allowbreak{}e\allowbreak{}/\allowbreak{}v\allowbreak{}4\allowbreak{}.\allowbreak{}5\allowbreak{}.\allowbreak{}0} in commit \texttt{6\allowbreak{}9\allowbreak{}2\allowbreak{}9\allowbreak{}1\allowbreak{}f\allowbreak{}a}. \texttt{d\allowbreak{}a\allowbreak{}g\allowbreak{}-\allowbreak{}e\allowbreak{}d\allowbreak{}g\allowbreak{}e\allowbreak{}s\allowbreak{}.\allowbreak{}t\allowbreak{}s} now routes same-column skip edges through a detour column beside intermediate nodes instead of drawing straight through the occluding node column, \texttt{G\allowbreak{}r\allowbreak{}i\allowbreak{}d\allowbreak{}S\allowbreak{}t\allowbreak{}a\allowbreak{}t\allowbreak{}e\allowbreak{}.\allowbreak{}a\allowbreak{}r\allowbreak{}r\allowbreak{}o\allowbreak{}w\allowbreak{}s} now tracks inbound multiplicity without collapsing routes into a boolean set, and the regression coverage in \texttt{d\allowbreak{}a\allowbreak{}g\allowbreak{}-\allowbreak{}e\allowbreak{}d\allowbreak{}g\allowbreak{}e\allowbreak{}s\allowbreak{}.\allowbreak{}t\allowbreak{}e\allowbreak{}s\allowbreak{}t\allowbreak{}.\allowbreak{}t\allowbreak{}s} / \texttt{d\allowbreak{}a\allowbreak{}g\allowbreak{}.\allowbreak{}t\allowbreak{}e\allowbreak{}s\allowbreak{}t\allowbreak{}.\allowbreak{}t\allowbreak{}s} locks the visual invariant down.

\section{Original Proposal}\label{record-289-original-proposal}

Legend: \href{https://github.com/flyingrobots/bijou/blob/1853ad902f6bf219326409e7f743a5a1d6517dbc/docs/method/legends/RE-runtime-engine.md}{RE --- Runtime Engine}

\section{Severity}\label{record-289-severity}

\textbf{Correctness bug, not cosmetic.} Downstream projects consume the DAG renderer to visualize graph-shaped data --- notably \texttt{g\allowbreak{}i\allowbreak{}t\allowbreak{}-\allowbreak{}w\allowbreak{}a\allowbreak{}r\allowbreak{}p}, a graph-database project that renders its graphs through bijou. For those consumers, the DAG renderer must \textbf{never hide or drop information}. A missing edge is a wrong answer, not an ugly one.

\section{Problem}\label{record-289-problem}

When a DAG edge skips a layer (source at layer L, destination at layer L+2 or beyond) and both endpoints end up in the same column, the vertical edge line is routed straight down through the intermediate layer --- passing through the rows where the intermediate layer's node box is drawn. The node box has higher render priority than the edge line, so the edge is visually erased.

The destination arrowhead survives because it lands in the gap row above the destination box. To a reader, it looks like the edge originates from the node directly above the destination, not from the real source.

Multiple edges terminating at the same destination cell also collide: \texttt{a\allowbreak{}r\allowbreak{}r\allowbreak{}o\allowbreak{}w\allowbreak{}s} is a \texttt{S\allowbreak{}e\allowbreak{}t\allowbreak{}<\allowbreak{}n\allowbreak{}u\allowbreak{}m\allowbreak{}b\allowbreak{}e\allowbreak{}r\allowbreak{}>} keyed by \texttt{e\allowbreak{}n\allowbreak{}c\allowbreak{}o\allowbreak{}d\allowbreak{}e\allowbreak{}A\allowbreak{}r\allowbreak{}r\allowbreak{}o\allowbreak{}w\allowbreak{}P\allowbreak{}o\allowbreak{}s\allowbreak{}(\allowbreak{}r\allowbreak{}o\allowbreak{}w\allowbreak{},\allowbreak{} \allowbreak{}c\allowbreak{}o\allowbreak{}l\allowbreak{})}, so only one arrowhead is drawn per destination even when the grid has two or more inbound edges.

\section{Evidence}\label{record-289-evidence}

Minimal repro (three nodes, edges \texttt{A\allowbreak{}→\allowbreak{}B}, \texttt{A\allowbreak{}→\allowbreak{}C}, \texttt{B\allowbreak{}→\allowbreak{}C}):

\begin{CodeBlock}
bijou_dag({ nodes: [
  { id: 'A', label: 'A', edges: ['B', 'C'] },
  { id: 'B', label: 'B', edges: ['C'] },
  { id: 'C', label: 'C', edges: [] },
]});
\end{CodeBlock}

Expected: two arrows into C (one from A, one from B). Actual: one arrow into C. The A→C edge is silently lost.

The loss was discovered while rendering the RE-017 byte-pipeline plan DAG --- \texttt{I\allowbreak{}-\allowbreak{}9\allowbreak{} \allowbreak{}→\allowbreak{} \allowbreak{}I\allowbreak{}I\allowbreak{}-\allowbreak{}1} and \texttt{I\allowbreak{}I\allowbreak{}I\allowbreak{}-\allowbreak{}5\allowbreak{} \allowbreak{}→\allowbreak{} \allowbreak{}I\allowbreak{}I\allowbreak{}I\allowbreak{}-\allowbreak{}6} (both layer-skipping, same-column) were invisible in the rendered output.

\section{Root Cause}\label{record-289-root-cause}

\begin{itemize}
\tightlist
\item
  \texttt{p\allowbreak{}a\allowbreak{}c\allowbreak{}k\allowbreak{}a\allowbreak{}g\allowbreak{}e\allowbreak{}s\allowbreak{}/\allowbreak{}b\allowbreak{}i\allowbreak{}j\allowbreak{}o\allowbreak{}u\allowbreak{}/\allowbreak{}s\allowbreak{}r\allowbreak{}c\allowbreak{}/\allowbreak{}c\allowbreak{}o\allowbreak{}r\allowbreak{}e\allowbreak{}/\allowbreak{}c\allowbreak{}o\allowbreak{}m\allowbreak{}p\allowbreak{}o\allowbreak{}n\allowbreak{}e\allowbreak{}n\allowbreak{}t\allowbreak{}s\allowbreak{}/\allowbreak{}d\allowbreak{}a\allowbreak{}g\allowbreak{}-\allowbreak{}e\allowbreak{}d\allowbreak{}g\allowbreak{}e\allowbreak{}s\allowbreak{}.\allowbreak{}t\allowbreak{}s\allowbreak{}:\allowbreak{}1\allowbreak{}4\allowbreak{}8\allowbreak{}-\allowbreak{}1\allowbreak{}4\allowbreak{}9} --- \texttt{m\allowbreak{}a\allowbreak{}r\allowbreak{}k\allowbreak{}E\allowbreak{}d\allowbreak{}g\allowbreak{}e} routes a same-column edge as a pure vertical segment from \texttt{s\allowbreak{}R\allowbreak{}o\allowbreak{}w} to \texttt{d\allowbreak{}R\allowbreak{}o\allowbreak{}w}, passing through intermediate layer rows.
\item
  \texttt{p\allowbreak{}a\allowbreak{}c\allowbreak{}k\allowbreak{}a\allowbreak{}g\allowbreak{}e\allowbreak{}s\allowbreak{}/\allowbreak{}b\allowbreak{}i\allowbreak{}j\allowbreak{}o\allowbreak{}u\allowbreak{}/\allowbreak{}s\allowbreak{}r\allowbreak{}c\allowbreak{}/\allowbreak{}c\allowbreak{}o\allowbreak{}r\allowbreak{}e\allowbreak{}/\allowbreak{}c\allowbreak{}o\allowbreak{}m\allowbreak{}p\allowbreak{}o\allowbreak{}n\allowbreak{}e\allowbreak{}n\allowbreak{}t\allowbreak{}s\allowbreak{}/\allowbreak{}d\allowbreak{}a\allowbreak{}g\allowbreak{}-\allowbreak{}r\allowbreak{}e\allowbreak{}n\allowbreak{}d\allowbreak{}e\allowbreak{}r\allowbreak{}.\allowbreak{}t\allowbreak{}s\allowbreak{}:\allowbreak{}3\allowbreak{}8\allowbreak{}3\allowbreak{}-\allowbreak{}4\allowbreak{}2\allowbreak{}0} --- \texttt{c\allowbreak{}e\allowbreak{}l\allowbreak{}l\allowbreak{}A\allowbreak{}t} resolves nodes before edges, so any grid cell inside a node box returns the box character, occluding the edge line underneath.
\item
  \texttt{p\allowbreak{}a\allowbreak{}c\allowbreak{}k\allowbreak{}a\allowbreak{}g\allowbreak{}e\allowbreak{}s\allowbreak{}/\allowbreak{}b\allowbreak{}i\allowbreak{}j\allowbreak{}o\allowbreak{}u\allowbreak{}/\allowbreak{}s\allowbreak{}r\allowbreak{}c\allowbreak{}/\allowbreak{}c\allowbreak{}o\allowbreak{}r\allowbreak{}e\allowbreak{}/\allowbreak{}c\allowbreak{}o\allowbreak{}m\allowbreak{}p\allowbreak{}o\allowbreak{}n\allowbreak{}e\allowbreak{}n\allowbreak{}t\allowbreak{}s\allowbreak{}/\allowbreak{}d\allowbreak{}a\allowbreak{}g\allowbreak{}-\allowbreak{}e\allowbreak{}d\allowbreak{}g\allowbreak{}e\allowbreak{}s\allowbreak{}.\allowbreak{}t\allowbreak{}s\allowbreak{}:\allowbreak{}1\allowbreak{}6\allowbreak{}2} --- \texttt{g\allowbreak{}.\allowbreak{}a\allowbreak{}r\allowbreak{}r\allowbreak{}o\allowbreak{}w\allowbreak{}s\allowbreak{}.\allowbreak{}a\allowbreak{}d\allowbreak{}d\allowbreak{}(\allowbreak{}e\allowbreak{}n\allowbreak{}c\allowbreak{}o\allowbreak{}d\allowbreak{}e\allowbreak{}A\allowbreak{}r\allowbreak{}r\allowbreak{}o\allowbreak{}w\allowbreak{}P\allowbreak{}o\allowbreak{}s\allowbreak{}(\allowbreak{}d\allowbreak{}R\allowbreak{}o\allowbreak{}w\allowbreak{},\allowbreak{} \allowbreak{}d\allowbreak{}s\allowbreak{}t\allowbreak{}C\allowbreak{})\allowbreak{})} collides for any two edges landing on the same destination row/col, discarding all but one arrowhead.
\end{itemize}

\section{Fix Options}\label{record-289-fix-options}

\begin{enumerate}
\def\labelenumi{\arabic{enumi}.}
\tightlist
\item
  \textbf{Detour around intermediate nodes.} When a layer-skipping edge would pass through a same-column intermediate node, route it through a jog column that's free of nodes at the intermediate layer. Biggest change; most visually correct.
\item
  \textbf{Column reassignment for skip edges.} When ordering columns, prefer to place layer-skipping edge endpoints in different columns so the straight-line route doesn't collide.
\item
  \textbf{Edge priority when no box char.} Keep the node-box priority for cells inside boxes, but detect pass-through and mark those cells with a ``node+edge'' flag so the box label area stays readable while still hinting that an edge crosses the box (e.g.~with a subdued glyph at the box border).
\item
  \textbf{Arrowhead multiplicity.} Change \texttt{a\allowbreak{}r\allowbreak{}r\allowbreak{}o\allowbreak{}w\allowbreak{}s} from \texttt{S\allowbreak{}e\allowbreak{}t\allowbreak{}<\allowbreak{}n\allowbreak{}u\allowbreak{}m\allowbreak{}b\allowbreak{}e\allowbreak{}r\allowbreak{}>} to \texttt{M\allowbreak{}a\allowbreak{}p\allowbreak{}<\allowbreak{}n\allowbreak{}u\allowbreak{}m\allowbreak{}b\allowbreak{}e\allowbreak{}r\allowbreak{},\allowbreak{} \allowbreak{}n\allowbreak{}u\allowbreak{}m\allowbreak{}b\allowbreak{}e\allowbreak{}r\allowbreak{}>} (count per cell) or allow the cell renderer to emit a divergent glyph when multiple edges converge.
\end{enumerate}

Option 1 is the most correct. Option 2 is probably the simplest and resolves the common case. At minimum, option 4 should ship so convergence isn't silently lost.

\section{Tests To Add}\label{record-289-tests-to-add}

\begin{itemize}
\tightlist
\item
  Three-node repro above, asserting \texttt{r\allowbreak{}e\allowbreak{}n\allowbreak{}d\allowbreak{}e\allowbreak{}r\allowbreak{}A\allowbreak{}c\allowbreak{}c\allowbreak{}e\allowbreak{}s\allowbreak{}s\allowbreak{}i\allowbreak{}b\allowbreak{}l\allowbreak{}e} shows \texttt{A\allowbreak{} \allowbreak{}→\allowbreak{} \allowbreak{}B\allowbreak{},\allowbreak{} \allowbreak{}C} and \texttt{B\allowbreak{} \allowbreak{}→\allowbreak{} \allowbreak{}C} (source of truth) and that the rendered output contains two arrowheads entering C's row range.
\item
  Four-node diamond (\texttt{A\allowbreak{} \allowbreak{}→\allowbreak{} \allowbreak{}B\allowbreak{},\allowbreak{} \allowbreak{}A\allowbreak{} \allowbreak{}→\allowbreak{} \allowbreak{}C\allowbreak{},\allowbreak{} \allowbreak{}B\allowbreak{} \allowbreak{}→\allowbreak{} \allowbreak{}D\allowbreak{},\allowbreak{} \allowbreak{}C\allowbreak{} \allowbreak{}→\allowbreak{} \allowbreak{}D}) --- asserting both inbound edges to D are visible.
\item
  A layer-skipping chain (\texttt{A\allowbreak{} \allowbreak{}→\allowbreak{} \allowbreak{}B\allowbreak{} \allowbreak{}→\allowbreak{} \allowbreak{}C}, plus \texttt{A\allowbreak{} \allowbreak{}→\allowbreak{} \allowbreak{}D} with D at layer 2) asserting A→D does not overlap B's column.
\end{itemize}

\section{Impact}\label{record-289-impact}

The pipe and accessible renderers (\texttt{r\allowbreak{}e\allowbreak{}n\allowbreak{}d\allowbreak{}e\allowbreak{}r\allowbreak{}P\allowbreak{}i\allowbreak{}p\allowbreak{}e}, \texttt{r\allowbreak{}e\allowbreak{}n\allowbreak{}d\allowbreak{}e\allowbreak{}r\allowbreak{}A\allowbreak{}c\allowbreak{}c\allowbreak{}e\allowbreak{}s\allowbreak{}s\allowbreak{}i\allowbreak{}b\allowbreak{}l\allowbreak{}e} in \texttt{d\allowbreak{}a\allowbreak{}g\allowbreak{}-\allowbreak{}r\allowbreak{}e\allowbreak{}n\allowbreak{}d\allowbreak{}e\allowbreak{}r\allowbreak{}.\allowbreak{}t\allowbreak{}s\allowbreak{}:\allowbreak{}4\allowbreak{}6\allowbreak{}9\allowbreak{}-\allowbreak{}5\allowbreak{}3\allowbreak{}0}) return the correct edge data, so any programmatic consumer is unaffected. Only the interactive (styled) renderer loses edges --- which is exactly the mode used for all human-facing output via the bijou MCP tool, the CLI, the bijou-tui panes, and the \texttt{g\allowbreak{}i\allowbreak{}t\allowbreak{}-\allowbreak{}w\allowbreak{}a\allowbreak{}r\allowbreak{}p} graph-database visualizations.

\section{Invariant To Enforce}\label{record-289-invariant-to-enforce}

\begin{quote}
The interactive DAG renderer MUST NOT hide, drop, or silently merge edges or arrowheads. Every edge in the input must be resolvable in the rendered output.
\end{quote}

This should become a repo invariant (\texttt{d\allowbreak{}o\allowbreak{}c\allowbreak{}s\allowbreak{}/\allowbreak{}i\allowbreak{}n\allowbreak{}v\allowbreak{}a\allowbreak{}r\allowbreak{}i\allowbreak{}a\allowbreak{}n\allowbreak{}t\allowbreak{}s\allowbreak{}/}) alongside a regression test that fails if any edge in a representative test graph is not visually present.

